\documentclass[12pt]{article}

\usepackage[utf8]{inputenc}
\usepackage[T1]{fontenc}
\usepackage{lmodern}
\usepackage[spanish]{babel}
\usepackage{booktabs}
\usepackage{amsmath}
\usepackage{forest}
\usepackage{float}
\usepackage{listings}
\usepackage{xcolor}
\usepackage{tikz}
\usetikzlibrary{arrows.meta, positioning, calc}


\definecolor{codegreen}{rgb}{0,0.6,0}
\definecolor{codegray}{rgb}{0.5,0.5,0.5}
\definecolor{codepurple}{rgb}{0.58,0,0.82}
\definecolor{backcolour}{rgb}{0.95,0.95,0.92}

% Definir colores para los nodos segun diferencia de valores
\definecolor{orangelight}{rgb}{1.0,0.8,0.6}
\definecolor{orangemedium}{rgb}{1.0,0.6,0.3}
\definecolor{orangedark}{rgb}{0.8,0.4,0.1}
\definecolor{bluelight}{rgb}{0.7,0.8,1.0}
\definecolor{bluemedium}{rgb}{0.4,0.6,1.0}
\definecolor{bluedark}{rgb}{0.2,0.4,0.8}

\lstdefinestyle{mystyle}{
    backgroundcolor=\color{backcolour},   
    commentstyle=\color{codegreen},
    keywordstyle=\color{magenta},
    numberstyle=\tiny\color{codegray},
    stringstyle=\color{codepurple},
    basicstyle=\ttfamily\footnotesize,
    breakatwhitespace=false,         
    breaklines=true,                 
    captionpos=b,                    
    keepspaces=true,                 
    numbers=left,                    
    numbersep=5pt,                  
    showspaces=false,                
    showstringspaces=false,
    showtabs=false,                  
    tabsize=2
}

\lstset{style=mystyle}

\sloppy
\setlength{\parindent}{0pt}

\begin{document}

% Título y materia
\begin{center}
  {\LARGE \textbf{Guía de Ejercicios\\Random Forest}}\\[0.5em]
  {Analítica de Datos, Universidad de San Andrés}
\end{center}

Si encuentran algún error en el documento o hay alguna duda, mandenmé un mail a rodriguezf@udesa.edu.ar y lo revisamos.

\section*{Ejercicios}

\begin{enumerate}
    \item ¿Cómo deben ser los árboles que componen un modelo de random forest? ¿Por qué? ¿Cómo hacemos para conseguirlo?

    \item Con este tipo de modelos ganamos poder predictivo pero perdemos interpretabilidad. ¿Qué podemos hacer en python para comprender la relevancia que tuvo cada variable al momento del entrenamiento?

    \item Teniendo en cuenta el siguiente modelo de random forest clasificar a las siguientes observaciones indicando la probabilidad.
    \begin{enumerate}
        \item Sex = 0; Class = 1; Embarked = A
        \item Sex = 1; Class = 2; Embarked = A
    \end{enumerate}

    \begin{center}
    \begin{tikzpicture}[
        scale=0.8, transform shape,
        every node/.style={align=center},
        treenode/.style={rectangle, draw, rounded corners, text width=8em, minimum height=3em, font=\scriptsize},
        line/.style={draw, -{Stealth[]}},
    ]
        % Tree 1 from image
        \node[treenode, fill=orangelight] (root) {Class <= 2.5\\gini = 0.469\\samples = 8\\value = [5, 3]\\class = No};

        \node[treenode, fill=bluemedium, below left=0.5cm and -1.5cm of root] (sex) {Sex <= 0.5\\gini = 0.375\\samples = 4\\value = [1, 3]\\class = Si};
        \node[treenode, fill=orangedark, below right=0.5cm and -1.5cm of root] (leaf1) {gini = 0.0\\samples = 4\\value = [4, 0]\\class = No};

        \node[treenode, fill=white, below left=0.5cm and -1cm of sex] (embarked) {Embarked\_B <= 0.5\\gini = 0.5\\samples = 2\\value = [1, 1]\\class = No};
        \node[treenode, fill=bluedark, below right=0.5cm and -1cm of sex] (leaf2) {gini = 0.0\\samples = 2\\value = [0, 2]\\class = Si};

        \node[treenode, fill=orangedark, below left=0.5cm and -1cm of embarked] (leaf3) {gini = 0.0\\samples = 1\\value = [1, 0]\\class = No};
        \node[treenode, fill=bluedark, below right=0.5cm and -1cm of embarked] (leaf4) {gini = 0.0\\samples = 1\\value = [0, 1]\\class = Si};

        \path[line] (root) -- node[midway, left] {True} (sex);
        \path[line] (root) -- node[midway, right] {False} (leaf1);

        \path[line] (sex) -- node[midway, left] {True} (embarked);
        \path[line] (sex) -- node[midway, right] {False} (leaf2);

        \path[line] (embarked) -- node[midway, left] {True} (leaf3);
        \path[line] (embarked) -- node[midway, right] {False} (leaf4);
    \end{tikzpicture}
    
    \vspace{2em}

    \begin{tikzpicture}[
        scale=0.8, transform shape,
        every node/.style={align=center},
        treenode/.style={rectangle, draw, rounded corners, text width=8em, minimum height=3em, font=\scriptsize},
        line/.style={draw, -{Stealth[]}},
    ]
        % Tree 2 from image
        \node[treenode, fill=white] (root) {Embarked\_A <= 0.5\\gini = 0.5\\samples = 8\\value = [4, 4]\\class = No};

        \node[treenode, fill=bluelight, below left=0.5cm and -1cm of root] (class) {Class <= 2.0\\gini = 0.444\\samples = 6\\value = [2, 4]\\class = Si};
        \node[treenode, fill=orangedark, below right=0.5cm and -1cm of root] (leaf1) {gini = 0.0\\samples = 2\\value = [2, 0]\\class = No};

        \node[treenode, fill=bluedark, below left=0.5cm and -1.5cm of class] (leaf2) {gini = 0.0\\samples = 1\\value = [0, 1]\\class = Si};
        \node[treenode, fill=bluelight, below right=0.5cm and -1.5cm of class] (sex) {Sex <= 0.5\\gini = 0.48\\samples = 5\\value = [2, 3]\\class = Si};

        \node[treenode, fill=orangedark, below left=0.5cm and -1cm of sex] (leaf3) {gini = 0.0\\samples = 1\\value = [1, 0]\\class = No};
        \node[treenode, fill=bluemedium, below right=0.5cm and -1cm of sex] (leaf4) {gini = 0.375\\samples = 4\\value = [1, 3]\\class = Si};

        \path[line] (root) -- node[midway, left] {True} (class);
        \path[line] (root) -- node[midway, right] {False} (leaf1);

        \path[line] (class) -- node[midway, left] {True} (leaf2);
        \path[line] (class) -- node[midway, right] {False} (sex);

        \path[line] (sex) -- node[midway, left] {True} (leaf3);
        \path[line] (sex) -- node[midway, right] {False} (leaf4);
    \end{tikzpicture}

    \vspace{2em}

    \begin{tikzpicture}[
        scale=0.8, transform shape,
        every node/.style={align=center},
        treenode/.style={rectangle, draw, rounded corners, text width=8em, minimum height=3em, font=\scriptsize},
        line/.style={draw, -{Stealth[]}},
    ]
      % Tree 3
      \node[treenode, fill=orangemedium] (root) {Embarked\_A <= 0.5\\gini = 0.375\\samples = 8\\value = [6, 2]\\class = Si};
      \node[treenode, fill=bluedark, below left=0.5cm and -1cm of root] (leaf1) {gini = 0.0\\samples = 2\\value = [0, 2]\\class = No};
      \node[treenode, fill=orangedark, below right=0.5cm and -1cm of root] (leaf2) {gini = 0.0\\samples = 6\\value = [6, 0]\\class = Si};
      
      \path [line] (root) -- node[midway, left] {True} (leaf1);
      \path [line] (root) -- node[midway, right] {False} (leaf2);
    \end{tikzpicture}
    \end{center}

    \item Entrenar un modelo de random forest para predecir si un animal es doméstico con el siguiente dataset de entrenamiento y teniendo en cuenta los siguientes hiperparámetros:
    \begin{itemize}
        \item Cantidad de árboles = 3
        \item Altura máxima de cada árbol = 2
        \item Cantidad de variables a tener en cuenta en cada partición = \(\sqrt{p}\)
    \end{itemize}

    \begin{table}[H]
        \centering
        \begin{tabular}{cccccc}
            \toprule
            \textbf{ID} & \textbf{tiene\_pelo} & \textbf{vuela} & \textbf{pone\_huevos} & \textbf{vive\_en\_casa} & \textbf{domestico} \\
            \midrule
            1 & sí & no & no & sí & sí \\
            2 & no & sí & sí & no & no \\
            3 & sí & no & sí & no & no \\
            4 & no & no & sí & no & no \\
            5 & sí & no & no & sí & sí \\
            6 & no & sí & sí & sí & sí \\
            7 & sí & no & no & no & no \\
            8 & no & no & sí & sí & sí \\
            9 & sí & no & sí & sí & sí \\
            10 & no & sí & sí & no & no \\
            \bottomrule
        \end{tabular}
    \end{table}

    Muestras bootstrap para cada árbol:
    \begin{itemize}
        \item Árbol 1: [6, 9, 4, 1, 5, 3, 3, 1, 8, 1]
        \item Árbol 2: [4, 1, 8, 5, 1, 7, 7, 7, 1, 5]
        \item Árbol 3: [3, 7, 1, 6, 6, 6, 1, 7, 3, 6]
    \end{itemize}
    La elección aleatoria de variables la pueden hacer en python.

\end{enumerate}

\newpage
\section*{Soluciones}

\begin{enumerate}
    \item Los árboles deben ser diversos y débiles individualmente. Se logra mediante bootstrap sampling (muestras con reemplazo), selección aleatoria de variables ($\sqrt{p}$ en cada split) y sin poda para permitir overfitting individual. Al promediar múltiples predictores diversos, se reduce la varianza total manteniendo bajo el sesgo.

    \item En Python usamos \texttt{rf.feature\_importances\_} (basado en reducción de impureza) o \texttt{permutation\_importance} (más robusta, mide degradación del modelo al permutar cada variable):
    
    \begin{lstlisting}[language=Python]
from sklearn.ensemble import RandomForestClassifier
from sklearn.inspection import permutation_importance

rf = RandomForestClassifier()
rf.fit(X_train, y_train)

# Metodo 1: Importancia por impureza
importancias = rf.feature_importances_

# Metodo 2: Permutation importance
perm_imp = permutation_importance(rf, X_test, y_test, n_repeats=10)
    \end{lstlisting}

    \item Variables codificadas: Sex (0=mujer, 1=hombre), Embarked\_A (1 si A, 0 si no), Embarked\_B (1 si B, 0 si no).
    
    \textbf{Observación 1 (Sex=0, Class=1, Embarked=A):}
    \begin{itemize}
        \item Árbol 1: Class$\leq$2.5 → Sex$\leq$0.5 → Embarked\_B$\leq$0.5 → No
        \item Árbol 2: Embarked\_A$\leq$0.5 (No) → No  
        \item Árbol 3: Embarked\_A$\leq$0.5 (No) → No
    \end{itemize}
    Predicción: No (probabilidad 1.0)
    
    \textbf{Observación 2 (Sex=1, Class=2, Embarked=A):}
    \begin{itemize}
        \item Árbol 1: Class$\leq$2.5 → Sex$\leq$0.5 (No) → Si
        \item Árbol 2: Embarked\_A$\leq$0.5 (No) → No
        \item Árbol 3: Embarked\_A$\leq$0.5 (No) → No
    \end{itemize}
    Predicción: No (prob. 0.67), Si (prob. 0.33)

    \item Dataset convertido a binario (sí=1, no=0). Con $\sqrt{4}=2$ variables por split y altura máxima 2, se construyen 3 árboles usando los bootstraps dados. La construcción manual requiere selección aleatoria de 2 variables en cada nodo. La idea es que hagamos nosotros el modelo de Random Forest sin usar el RandomForestClassifier de sklearn.
    
    \begin{lstlisting}[language=Python]
import numpy as np
import pandas as pd
from sklearn.tree import DecisionTreeClassifier

# Crear dataset
data = {
    'pelo': [1,0,1,0,1,0,1,0,1,0],
    'vuela': [0,1,0,0,0,1,0,0,0,1], 
    'huevos': [0,1,1,1,0,1,0,1,1,1],
    'casa': [1,0,0,0,1,1,0,1,1,0],
    'domestico': [1,0,0,0,1,1,0,1,1,0]
}

df = pd.DataFrame(data)
X = df[['pelo', 'vuela', 'huevos', 'casa']]
y = df['domestico']

# Bootstraps dados
bootstrap1 = [5,8,3,0,4,2,2,0,7,0]  # IDs - 1 (0-indexed)
bootstrap2 = [3,0,7,4,0,6,6,6,0,4]
bootstrap3 = [2,6,0,5,5,5,0,6,2,5]

# Entrenar arboles individuales
trees = []
for i, bootstrap in enumerate([bootstrap1, bootstrap2, bootstrap3]):
    X_boot = X.iloc[bootstrap]
    y_boot = y.iloc[bootstrap]
    
    tree = DecisionTreeClassifier(
        max_depth=2,
        max_features=2,  # sqrt(4) = 2
        random_state=i
    )
    tree.fit(X_boot, y_boot)
    trees.append(tree)

# Hacer predicciones con el Random Forest
def predict_random_forest(X_new, trees):
    predicciones = []
    for tree in trees:
        pred = tree.predict(X_new)[0]  # Solo una observacion
        predicciones.append(pred)
    
    # Votar por mayoria simple
    votos_si = sum(predicciones)
    votos_no = len(predicciones) - votos_si
    
    if votos_si > votos_no:
        return 1  # domestico
    else:
        return 0  # no domestico

# Ejemplo de prediccion
nueva_obs = [[1, 0, 0, 1]]  # pelo=si, vuela=no, huevos=no, casa=si
prediccion = predict_random_forest(nueva_obs, trees)
print(f"Prediccion: {'domestico' if prediccion == 1 else 'no domestico'}")
    \end{lstlisting}

\end{enumerate}

\end{document}
